% ==========================================================
% Requiere: amsmath, amssymb, amsthm, enumitem
% Entornos usados: lema, theorem, corollary, remark, proof
% Notación: se usa ^{T} para la traspuesta y ^{*} para la adjunta.
% ==========================================================

\chapter{Schur y bases ortonormales de autovectores}\label{cap:schur}

El Capítulo~\ref{cap:autovalores} terminó con una respuesta y una frustración. La
respuesta: una matriz es diagonalizable exactamente cuando tiene una base de autovectores
independientes, y entonces $A=X\Lambda X^{-1}$ convierte potencias, polinomios y sistemas
dinámicos en cuentas con números. La frustración: hay matrices que no la tienen, y no
parece haber nada que hacer con ellas.

Este capítulo muestra que sí lo hay, y que el remedio pasa por bajar una pretensión y
subir otra. Bajamos la pretensión de llegar hasta una matriz \emph{diagonal} y nos
conformamos con una \emph{triangular}; a cambio, subimos la exigencia sobre la base y la
pedimos \textbf{ortonormal}. El resultado es el teorema de Schur, y lo notable es que
vale \emph{siempre}, sin ninguna hipótesis: toda matriz cuadrada, diagonalizable o no,
real o compleja, se escribe como
\[
A=UTU^{*},\qquad U \text{ unitaria},\quad T \text{ triangular superior}.
\]
Cambiar $X^{-1}$ por $U^{*}$ no es un detalle: las matrices unitarias son las que
conservan normas y ángulos --- lo probamos en la Sección~\ref{sec:notacion-compleja}
---, así que una factorización con $U$ es estable y una con $X$ cualquiera puede no
serlo. Por eso esta es la factorización que usan de verdad los algoritmos numéricos.

Con Schur en la mano, dos preguntas se vuelven fáciles. La primera es qué información se
lee en la diagonal de $T$: allí están los autovalores con sus multiplicidades, y de ahí
salen en tres renglones las fórmulas de la traza y el determinante que el
Capítulo~\ref{cap:autovalores} había dejado anunciadas. La segunda es cuándo se puede
llegar hasta una matriz diagonal manteniendo la base ortonormal, es decir, cuándo
\emph{sí} hay una base ortonormal de autovectores. La respuesta es sorprendentemente
limpia:
\[
A \text{ tiene una base ortonormal de autovectores}
\iff
AA^{*}=A^{*}A ,
\]
una condición que se verifica multiplicando dos matrices. Las matrices que la cumplen se
llaman \textbf{normales}, e incluyen a las simétricas, las hermíticas, las ortogonales,
las unitarias y las de permutación: casi todas las matrices con nombre propio que fuimos
encontrando. El caso real y simétrico de este resultado es el \emph{teorema espectral},
$A=Q\Lambda Q^{T}$, que es la herramienta con la que se construyen las matrices definidas
positivas del Capítulo~\ref{cap:sdp} y la SVD del Capítulo~\ref{cap:SVD}.

Cerramos con una consecuencia cuantitativa: para una matriz normal, el tamaño de la
matriz --- medido con la norma de Frobenius de la Sección~\ref{sec:frobenius} --- está
íntegramente repartido entre sus autovalores. Eso es lo que permite decir cuánto se
pierde al quedarse sólo con los más grandes, que es la idea que sostiene la aproximación
de rango bajo y el análisis de componentes principales.

\section{El caso complejo: notación y matrices unitarias}
\label{sec:notacion-compleja}

Como en el Capítulo~\ref{cap:autovalores}, acá no alcanza con $\R$: los autovalores
viven en $\mathbb{C}$ y los autovectores en $\mathbb{C}^n$. Antes de empezar fijamos la
notación y probamos las pocas propiedades que vamos a usar todo el tiempo. Nada de esto
es difícil; lo único que hay que hacer es acordarse de conjugar.

\subsection*{La adjunta \texorpdfstring{$A^{*}$}{A*}}

En la Sección~\ref{sec:producto-interno-complejo} vimos que el producto interno de
$\mathbb{C}^n$ es $\overline{\bx}^{T}\by$, con el primer factor conjugado. La operación
``conjugar y trasponer'' aparece tantas veces que conviene darle nombre y símbolo.

\begin{definicion}[Adjunta]\label{def:adjunta}
Sea $A\in\mathbb{C}^{m\times n}$. La \textbf{adjunta} (o traspuesta conjugada) de $A$ es
la matriz $A^{*}\in\mathbb{C}^{n\times m}$ dada por
\[
\boxed{\ A^{*}=\overline{A}^{\,T},\qquad (A^{*})_{ij}=\overline{a_{ji}}\ }
\]
\end{definicion}

Con esta notación el producto interno de $\mathbb{C}^n$ se escribe simplemente
$\bx^{*}\by$: un vector columna es una matriz de $n\times1$, y su adjunta es el vector
fila conjugado. Las reglas son las de la traspuesta, con un conjugado encima:
\[
(A^{*})^{*}=A,
\qquad
(A+B)^{*}=A^{*}+B^{*},
\qquad
(cA)^{*}=\overline{c}\,A^{*},
\qquad
\boxed{\ (AB)^{*}=B^{*}A^{*}\ }
\]
La última, que vuelve a invertir el orden, se verifica igual que para la traspuesta:
$\overline{(AB)}^{\,T}=(\overline{A}\,\overline{B})^{T}
=\overline{B}^{\,T}\overline{A}^{\,T}=B^{*}A^{*}$.

El nombre \emph{adjunta} viene de la propiedad siguiente, que es la razón de ser de toda
esta notación: $A^{*}$ es la matriz que permite pasar $A$ de un lado al otro del
producto interno.

\begin{proposicion}\label{prop:adjunta}
Para todos $\bx\in\mathbb{C}^{n}$, $\by\in\mathbb{C}^{m}$ y $A\in\mathbb{C}^{m\times n}$,
\[
(A\bx)^{*}\by=\bx^{*}(A^{*}\by) .
\]
\end{proposicion}

\begin{proof}
$(A\bx)^{*}\by=(\bx^{*}A^{*})\by=\bx^{*}(A^{*}\by)$, usando $(AB)^{*}=B^{*}A^{*}$ y la
asociatividad del producto.
\end{proof}

\subsection*{Hermíticas y unitarias}

Dos familias de matrices complejas llevan nombre propio, y son las versiones complejas
de dos que ya conocemos.

\begin{definicion}[Hermítica y unitaria]\label{def:hermitica-unitaria}
Sea $A\in\mathbb{C}^{n\times n}$. Decimos que
\begin{itemize}[itemsep=2pt,topsep=3pt]
\item $A$ es \textbf{hermítica} si $A^{*}=A$, es decir, si $a_{ij}=\overline{a_{ji}}$;
\item $A$ es \textbf{unitaria} si $A^{*}A=I$, es decir, si sus columnas forman una base
ortonormal de $\mathbb{C}^n$ con el producto interno complejo.
\end{itemize}
\end{definicion}

El diccionario con el caso real es inmediato: si $A$ tiene todas sus entradas reales
entonces $\overline{A}=A$ y por lo tanto $A^{*}=A^{T}$, de modo que
\[
\text{hermítica}\ \longleftrightarrow\ \text{simétrica},
\qquad
\text{unitaria}\ \longleftrightarrow\ \text{ortogonal}.
\]
Todo lo que digamos de hermíticas y unitarias contiene, como caso particular, lo que ya
sabemos de simétricas y ortogonales. Observemos de paso que una matriz hermítica tiene
la diagonal real, porque $a_{ii}=\overline{a_{ii}}$.

\begin{observacion}[Cuidado con la palabra]
En el Capítulo~\ref{cap:LU} hablamos de matrices triangulares ``de diagonal unitaria'',
queriendo decir que los elementos de la diagonal valen $1$. Es un uso distinto y
desafortunado de la misma palabra: una matriz \emph{unitaria} en el sentido de la
definición anterior no tiene por qué tener unos en la diagonal, ni al revés. El contexto
siempre deja claro cuál de los dos sentidos está en juego.
\end{observacion}

Lo único que necesitamos de las unitarias es que no deforman: ésta es la propiedad que
hace que la factorización de Schur sea numéricamente estable, y es la versión compleja
de la Proposición~\ref{prop:preserva} del Capítulo~\ref{cap:ortogonalidad}, que allí
probamos para matrices reales.

\begin{proposicion}[Las unitarias preservan normas y ángulos]\label{prop:unitarias-preservan}
Sea $U\in\mathbb{C}^{n\times n}$ unitaria. Entonces, para todos
$\bx,\by\in\mathbb{C}^{n}$,
\[
(U\bx)^{*}(U\by)=\bx^{*}\by
\qquad\text{y}\qquad
\|U\bx\|=\|\bx\| .
\]
Además $U$ es invertible con $U^{-1}=U^{*}$, y $U^{*}$ es también unitaria.
\end{proposicion}

\begin{proof}
La primera igualdad es una cuenta:
\[
(U\bx)^{*}(U\by)=\bx^{*}U^{*}U\by=\bx^{*}I\by=\bx^{*}\by .
\]
Tomando $\by=\bx$ queda $\|U\bx\|^{2}=\|\bx\|^{2}$, y como las normas son no negativas,
$\|U\bx\|=\|\bx\|$. Para la invertibilidad: $U$ es cuadrada y $U^{*}U=I$ dice que
$U^{*}$ es inversa a izquierda; para una matriz cuadrada eso basta para que sea la
inversa, de modo que $U^{-1}=U^{*}$ y también $UU^{*}=I$. Esta última igualdad, leída
como $(U^{*})^{*}U^{*}=I$, dice que $U^{*}$ es unitaria.
\end{proof}

Conservar el producto interno es conservar normas \emph{y} ángulos: si dos vectores eran
ortogonales, sus imágenes lo siguen siendo. Por eso multiplicar por una matriz unitaria
no agranda errores de redondeo, y por eso todas las factorizaciones que usan los
algoritmos numéricos --- $QR$, Schur, SVD --- están hechas con matrices de esta clase.

\section{El teorema de Schur}\label{sec:teo-schur}

\begin{theorem}[Schur]\label{teo:schur}   % <-- CORREGIDO: es teorema, no lema
Para toda $A\in\mathbb{C}^{n\times n}$ existe una matriz unitaria $U$ tal que
\[
T:=U^{*}AU \quad\text{es triangular superior,}\qquad\text{equivalentemente}\qquad A=UTU^{*}.
\]
Además, la diagonal de $T$ está formada por los autovalores de $A$ (contados con
multiplicidad). La factorización $A=UTU^{*}$ se llama \emph{descomposición de Schur} de $A$.
\end{theorem}

\begin{proof}[Demostración por inducción en $n$]
\textbf{Caso base $n=1$.} Trivial.

\smallskip
\textbf{Paso inductivo.} Supongamos el enunciado cierto para tamaño $n-1$ y sea
$A\in\mathbb{C}^{n\times n}$. Por el teorema fundamental del álgebra, el polinomio
característico de $A$ tiene una raíz $\lambda\in\mathbb{C}$; sea $v\neq 0$ un autovector
asociado. Normalizamos $\|v\|=1$ y extendemos $\{v\}$ a una base ortonormal de
$\mathbb{C}^n$, $\{v,v_2,\dots,v_n\}$ (por ejemplo con Gram--Schmidt), y sea
\[
U_1=[\,v\;\;v_2\;\cdots\;v_n\,].
\]
Entonces $U_1$ es unitaria y, escribiendo $B:=U_1^{*}AU_1$, se tiene
\[
B e_1 \;=\; U_1^{*}A\,U_1e_1\;=\;U_1^{*}A v \;=\; U_1^{*}(\lambda v)\;=\;\lambda\, U_1^{*}v\;=\;\lambda e_1 ,
\]
con $e_1=(1,0,\dots,0)^{T}$. Por lo tanto la primera columna de $B$ es
$(\lambda,0,\dots,0)^{T}$ y $B$ adopta la forma por bloques
\[
B=\begin{bmatrix}
\lambda & *\\[2pt]
0 & A_1
\end{bmatrix},
\qquad A_1\in\mathbb{C}^{(n-1)\times(n-1)}.
\]
Por hipótesis inductiva existe una unitaria $U_2\in\mathbb{C}^{(n-1)\times(n-1)}$
tal que $U_2^{*}A_1U_2=T_1$ es triangular superior. Definimos
\[
U \;=\; U_1\begin{bmatrix}1&0\\[2pt]0&U_2\end{bmatrix},
\]
que es unitaria por ser producto de unitarias, y calculamos
\[
U^{*}AU
=\begin{bmatrix}1&0\\[2pt]0&U_2^{*}\end{bmatrix}
\begin{bmatrix}\lambda & *\\[2pt]0 & A_1\end{bmatrix}
\begin{bmatrix}1&0\\[2pt]0&U_2\end{bmatrix}
=
\begin{bmatrix}\lambda & *\\[2pt]0 & T_1\end{bmatrix},
\]
que es triangular superior. La diagonal está formada por $\lambda$ seguido de la
diagonal de $T_1$, que por inducción son los autovalores de $A_1$ y, por ende, de $A$
(pues $A$ y $B$ son semejantes y $\det(B-zI)=(\lambda-z)\det(A_1-zI)$).
Esto cierra la inducción.
\end{proof}

\begin{corollary}[Las multiplicidades se leen en la diagonal]\label{cor:mult-diagonal}
Con las notaciones del Teorema~\ref{teo:schur}, el polinomio característico de $A$ es
\[
\det(A-zI)=\det(T-zI)=\prod_{k=1}^{n}(t_{kk}-z),
\]
de modo que cada autovalor aparece en la diagonal de $T$ tantas veces como indica su
\emph{multiplicidad algebraica}. En consecuencia, si $\lambda$ tiene multiplicidad
algebraica $m$, entonces $T-\lambda I$ tiene exactamente $m$ ceros en la diagonal,
$\rg(T-\lambda I)\ge n-m$ y por lo tanto
\[
\dim E_\lambda=\dim N(A-\lambda I)\le m :
\]
\emph{la multiplicidad geométrica nunca supera a la algebraica}.
\end{corollary}

\begin{proof}
$A-zI=U(T-zI)U^{*}$ es semejante a $T-zI$, y el determinante de una triangular es el
producto de su diagonal. Para la segunda parte, $A-\lambda I$ y $T-\lambda I$ son
semejantes, así que tienen el mismo rango y espacios nulos de la misma dimensión; y una
matriz triangular con $n-m$ elementos diagonales no nulos tiene al menos $n-m$ columnas
independientes.
\end{proof}

\begin{remark}
La desigualdad entre multiplicidades ya la habíamos demostrado, por otro camino, en el
Ejercicio de desarrollo~\ref{ent:multiplicidades} del
Capítulo~\ref{cap:autovalores}. Vale la pena comparar las dos pruebas: allá se extendía
una base del autoespacio y se miraba un bloque; acá se triangula de una vez y la
desigualdad se lee en la diagonal.
\end{remark}

\begin{remark}[El orden de la diagonal se puede elegir]  % <-- NUEVO
En el paso inductivo, $\lambda$ es \emph{cualquier} autovalor de $A$. Repitiendo el argumento
se ve que, dado un orden prefijado $\lambda_1,\dots,\lambda_n$ de los autovalores de $A$
(con multiplicidad), existe una unitaria $U$ tal que $U^{*}AU$ es triangular superior
con esa diagonal en ese orden. Esto se usa, por ejemplo, para poner primero el autovalor
de mayor módulo.
\end{remark}

\begin{corollary}[Traza, determinante, potencias y polinomios]
\label{cor:traza-det}
Sea $A\in\mathbb{C}^{n\times n}$ con autovalores $\lambda_1,\dots,\lambda_n$ contados
con multiplicidad algebraica. Entonces
\[
\tr(A)=\lambda_1+\cdots+\lambda_n,
\qquad
\det(A)=\lambda_1\lambda_2\cdots\lambda_n .
\]
Más aún, para todo $m\in\mathbb{N}$ los autovalores de $A^{m}$ son
$\lambda_1^{m},\dots,\lambda_n^{m}$, y para todo polinomio $p$ los autovalores de
$p(A)$ son $p(\lambda_1),\dots,p(\lambda_n)$, en todos los casos contados con
multiplicidad. Si $A$ es invertible, los autovalores de $A^{-1}$ son
$\lambda_1^{-1},\dots,\lambda_n^{-1}$.
\end{corollary}

\begin{proof}
Escribamos $A=UTU^{*}$ con $U$ unitaria y $T$ triangular superior cuya diagonal
es $\lambda_1,\dots,\lambda_n$ (Teorema~\ref{teo:schur}). Como
$\tr(XY)=\tr(YX)$,
\[
\tr(A)=\tr(UTU^{*})=\tr(TU^{*}U)
=\tr(T)=\lambda_1+\cdots+\lambda_n ,
\]
y como el determinante es multiplicativo y $\det(U^{*})=\det(U)^{-1}$,
\[
\det(A)=\det(U)\det(T)\det(U^{*})=\det(T)=\lambda_1\cdots\lambda_n ,
\]
porque el determinante de una matriz triangular es el producto de su diagonal.

Para las potencias, $A^{m}=UT^{m}U^{*}$, y el producto de matrices triangulares
superiores es triangular superior con la diagonal multiplicada entrada a entrada:
la diagonal de $T^{m}$ es $\lambda_1^{m},\dots,\lambda_n^{m}$. Como $A^{m}$ es
semejante a $T^{m}$, esos son sus autovalores. El mismo argumento, aplicado a
$p(T)$, da el caso de un polinomio; y si $A$ es invertible, ningún $\lambda_i$ es
cero, $T$ es invertible y la diagonal de $T^{-1}$ es $\lambda_i^{-1}$.
\end{proof}

Una consecuencia pequeña pero que vamos a usar varias veces cierra una deuda
del Capítulo~\ref{cap:ortogonalidad}.

\begin{corollary}[La traza de una proyección es su rango]\label{cor:traza-proyeccion}
Sea $P\in\R^{n\times n}$ con $P^2=P$ (una proyección, ortogonal o no).
Entonces los autovalores de $P$ son $0$ y $1$, y
\[
\boxed{\ \rg(P)=\tr(P)\ }
\]
\end{corollary}

\begin{proof}
Si $Px=\lambda x$ con $x\neq0$, entonces
$\lambda x = Px = P^2x = \lambda Px=\lambda^2x$, de donde
$\lambda^2=\lambda$ y $\lambda\in\{0,1\}$. El espacio propio de $\lambda=1$ es
$C(P)$: si $Px=x$ entonces $x\in C(P)$, y recíprocamente si $x=Py$ entonces
$Px=P^2y=Py=x$. El de $\lambda=0$ es $N(P)$. Como
$\R^n=C(P)\oplus N(P)$ ---todo $x$ se escribe $x=Px+(x-Px)$ y la
intersección es nula---, la suma de las dimensiones de los dos espacios
propios es $n$: $P$ es diagonalizable y sus autovalores son $1$ repetido
$\rg(P)$ veces y $0$ repetido $n-\rg(P)$ veces. Por el
Corolario~\ref{cor:traza-det}, $\tr(P)$ es la suma de esos autovalores, es
decir $\rg(P)$.
\end{proof}

\begin{remark}
Estas fórmulas se anunciaron en el Capítulo~\ref{cap:autovalores} sin demostración.
Vale la pena notar lo que aportó Schur: allí sabíamos que \emph{cada} $\lambda_i^{m}$
era un autovalor de $A^{m}$, pero no que fueran \emph{todos} ni con qué multiplicidad.
Triangular de una sola vez resuelve las dos cosas.
\end{remark}

\begin{remark}[Qué dice y qué no dice Schur]  % <-- NUEVO
Schur vale para \emph{toda} matriz compleja, sin ninguna hipótesis: siempre hay
una base ortonormal en la que $A$ es triangular. Lo que \emph{no} dice es que esa base
sea de autovectores: solo la primera columna de $U$ lo es. La pregunta natural
es entonces \emph{cuándo} se puede llegar hasta una matriz diagonal. La respuesta es
el Teorema~\ref{teo:normal-iff-ortonormal}.
\end{remark}

\begin{example}[Una descomposición de Schur a mano]\label{ej:schur-numerico}
Sea
\[
A=\begin{bmatrix}1&1\\ 0&2\end{bmatrix}.
\]
Sus autovalores son $1$ y $2$ (es triangular). Tomemos el autovalor $2$ para que la
primera columna de $U$ sea un autovector suyo: de $(A-2I)\n{v}=\n{0}$ sale
$\n{v}=(1,1)^{T}$, que normalizado es $\n{q}_1=\tfrac1{\sqrt2}(1,1)^{T}$. Completamos con
$\n{q}_2=\tfrac1{\sqrt2}(1,-1)^{T}$, ortogonal al anterior, y formamos
\[
U=\frac{1}{\sqrt2}\begin{bmatrix}1&1\\ 1&-1\end{bmatrix}
\qquad\Longrightarrow\qquad
T=U^{*}AU=\begin{bmatrix}2&-1\\ 0&1\end{bmatrix}.
\]
Se ve lo que anuncia la observación anterior: la diagonal de $T$ tiene los autovalores en
el orden elegido, la primera columna de $U$ es un autovector, y la segunda \emph{no} lo es
--- si lo fuera, $T$ sería diagonal. Y no puede serlo: $A$ no es normal
($AA^{T}\neq A^{T}A$), aunque sí es diagonalizable, con autovectores $(1,0)^T$ y
$(1,1)^T$ que son independientes pero no ortogonales. \hfill$\square$
\end{example}

\section{Matrices normales}\label{sec:normales}

\paragraph{Definición (matriz normal).}
Una matriz \(A\in\mathbb{C}^{n\times n}\) es \emph{normal} si
\[
A A^{*}=A^{*}A,
\]
donde \(A^{*}=\overline{A}^{T}\) denota la adjunta (conjugada traspuesta) de \(A\).

\medskip
\begin{lema}[Triangular normal \(\Rightarrow\) diagonal]\label{lema:triangular-normal}
Si \(T\in\mathbb{C}^{n\times n}\) es triangular superior y normal, entonces \(T\) es diagonal.
\end{lema}

\begin{proof}[Demostración por inducción en $n$]
\textbf{Caso base $n=1$.} Trivial.

\medskip
\textbf{Paso inductivo.} Supongamos el enunciado cierto para matrices
$(n-1)\times(n-1)$. Sea $T\in\mathbb{C}^{n\times n}$ triangular superior y normal.
Escribimos $T$ en bloques
\[
T=\begin{bmatrix}
\lambda & w^{*}\\
0 & S
\end{bmatrix},
\qquad
T^{*}=\begin{bmatrix}
\overline{\lambda} & 0\\
w & S^{*}
\end{bmatrix},
\]
donde $\lambda\in\mathbb{C}$, $w\in\mathbb{C}^{\,n-1}$ y $S\in\mathbb{C}^{(n-1)\times(n-1)}$
es triangular superior. Calculamos
\[
T^{*}T=
\begin{bmatrix}
|\lambda|^{2} & \overline{\lambda}\,w^{*}\\[2pt]
\lambda\,w & ww^{*}+S^{*}S
\end{bmatrix},
\qquad
TT^{*}=
\begin{bmatrix}
|\lambda|^{2}+\|w\|_{2}^{2} & w^{*}S^{*}\\[2pt]
Sw & SS^{*}
\end{bmatrix}.
\]
Como $T$ es normal, $T^{*}T=TT^{*}$. Comparando el bloque $(1,1)$ se obtiene
$|\lambda|^{2}=|\lambda|^{2}+\|w\|_{2}^{2}$, luego $\|w\|_{2}^{2}=0$ y por tanto $w=0$.
Con $w=0$, la igualdad de los bloques $(2,2)$ da $S^{*}S=SS^{*}$, o sea,
$S$ es normal. Como además $S$ es triangular superior, por hipótesis inductiva
$S$ es diagonal. De $w=0$ se sigue que
\[
T=\begin{bmatrix}\lambda&0\\[2pt]0&S\end{bmatrix}
\]
es diagonal. Esto completa la inducción.
\end{proof}

% ----------------------------------------------------------
\begin{lema}[Propiedades básicas de una normal]\label{lema:normal-basico}  % <-- NUEVO
Sea $A\in\mathbb{C}^{n\times n}$ normal. Entonces:
\begin{enumerate}[label=(\roman*)]
\item $\|Av\|=\|A^{*}v\|$ para todo $v\in\mathbb{C}^{n}$.
\item $Av=\lambda v \iff A^{*}v=\overline{\lambda}\,v$. Es decir, $A$ y $A^{*}$ tienen
      los mismos autovectores, con autovalores conjugados.
\item Autovectores asociados a autovalores \emph{distintos} son ortogonales.
\end{enumerate}
\end{lema}

\begin{proof}
(i) $\|Av\|^{2}=v^{*}A^{*}Av=v^{*}AA^{*}v=\|A^{*}v\|^{2}$.

\smallskip
(ii) Si $A$ es normal, también lo es $A-\lambda I$, porque
$(A-\lambda I)(A-\lambda I)^{*}=AA^{*}-\lambda A^{*}-\overline{\lambda}A+|\lambda|^{2}I$
coincide con $(A-\lambda I)^{*}(A-\lambda I)$ usando $AA^{*}=A^{*}A$.
Aplicando (i) a $A-\lambda I$:
\[
\|(A-\lambda I)v\|=\|(A-\lambda I)^{*}v\|=\|(A^{*}-\overline{\lambda}I)v\| ,
\]
y un lado se anula si y solo si se anula el otro.

\smallskip
(iii) Sean $Av=\lambda v$ y $Aw=\mu w$ con $\lambda\neq\mu$. Por (ii),
$A^{*}w=\overline{\mu}\,w$. Entonces
\[
\lambda\,(w^{*}v)=w^{*}(Av)=(A^{*}w)^{*}v=(\overline{\mu}\,w)^{*}v=\mu\,(w^{*}v),
\]
luego $(\lambda-\mu)\,w^{*}v=0$ y, como $\lambda\neq\mu$, resulta $w^{*}v=0$.
\end{proof}

\begin{theorem}[Caracterización vía Schur]\label{teo:normal-iff-ortonormal}
Sea \(A\in\mathbb{C}^{n\times n}\). Son equivalentes:
\begin{enumerate}
\item \(A\) es normal, es decir, \(AA^{*}=A^{*}A\).
\item \(A\) admite una base \emph{ortonormal} de autovectores; equivalentemente, existen
      $U$ unitaria y $D$ diagonal tales que \(A=U D U^{*}\)
      (\emph{$A$ es unitariamente diagonalizable}).   % <-- CORREGIDO: "unitariamente"
\end{enumerate}
\end{theorem}

\begin{proof}
\((2)\Rightarrow(1)\): Si \(A=U D U^{*}\) con \(U\) unitaria y \(D\) diagonal, entonces
$A^{*}=(UDU^{*})^{*}=UD^{*}U^{*}$, y como $D$ es diagonal $D^{*}=\overline{D}$
(la diagonal conjugada). Por lo tanto   % <-- CORREGIDO: aparece D^*, no \overline D de la nada
\[
AA^{*}=U D U^{*}\,U \overline{D} U^{*}=U (D\overline{D}) U^{*}
      =U (\overline{D}D) U^{*}=U\overline{D}U^{*}\,UDU^{*}=A^{*}A,
\]
porque dos matrices diagonales siempre conmutan.

\smallskip
\((1)\Rightarrow(2)\): Tomamos la descomposición de Schur \(A=U T U^{*}\) con \(T\)
triangular superior (Teorema~\ref{teo:schur}). Entonces
\[
AA^{*}=U T U^{*}\, U T^{*} U^{*}=U(TT^{*})U^{*},\qquad
A^{*}A=U (T^{*} T) U^{*}.
\]
Como $U$ es invertible, la igualdad \(AA^{*}=A^{*}A\) equivale a \(TT^{*}=T^{*}T\),
o sea, \(T\) es normal. Por el Lema~\ref{lema:triangular-normal}, \(T\) es diagonal.
Luego $A=UDU^{*}$ con $D=T$ diagonal, y las columnas de $U$ forman una base ortonormal
de autovectores de $A$.
\end{proof}

\begin{corollary}[Descomposición espectral]\label{cor:espectral}   % <-- NUEVO
Si $A$ es normal, con $U=[q_1\;\cdots\;q_n]$ y $D=\operatorname{diag}(\lambda_1,\dots,\lambda_n)$
como arriba, entonces
\[
A=\sum_{k=1}^{n}\lambda_k\, q_k q_k^{*},
\qquad\text{y además}\qquad
I=\sum_{k=1}^{n} q_k q_k^{*}.
\]
Cada $q_kq_k^{*}$ es el proyector ortogonal sobre la recta generada por $q_k$:
la matriz $A$ queda descompuesta como una suma de proyecciones ortogonales
\emph{pesadas por los autovalores}. En consecuencia, para todo $m\in\mathbb{N}$,
\[
A^{m}=\sum_{k=1}^{n}\lambda_k^{\,m}\, q_k q_k^{*},
\]
y más generalmente $p(A)=\sum_k p(\lambda_k)\,q_kq_k^{*}$ para cualquier polinomio $p$.
\end{corollary}

\begin{proof}
$A=UDU^{*}$ y el producto $UDU^{*}$, escrito por columnas, es $\sum_k \lambda_k q_kq_k^{*}$;
con $D=I$ se obtiene la segunda identidad. Para las potencias,
$A^{m}=UD^{m}U^{*}$ pues $U^{*}U=I$.
\end{proof}

\begin{remark}[Lectura para datos]   % <-- NUEVO
El Corolario~\ref{cor:espectral} es el motivo por el cual los autovalores ``mandan''
en las aplicaciones: aplicar $A$ repetidas veces amplifica las direcciones $q_k$ con
$|\lambda_k|$ grande y aplasta las de $|\lambda_k|$ chico. De ahí salen el método
de la potencia, la convergencia de cadenas de Markov, los filtros de señales y
la aproximación de rango bajo.
\end{remark}

\begin{entrega}{Matrices que conmutan y diagonalización simultánea}
\label{ent:conmutan}
El teorema espectral diagonaliza una matriz normal. Este ejercicio responde una pregunta
natural: ¿cuándo se pueden diagonalizar \emph{dos} matrices con la \emph{misma} base?

\begin{enumerate}[label=\textbf{(\alph*)},itemsep=3pt,topsep=4pt]
\item Sean $A,B\in\mathbb{C}^{n\times n}$ con $AB=BA$. Probar que cada autoespacio de
$A$ es invariante por $B$: si $Av=\lambda v$, entonces $A(Bv)=\lambda(Bv)$.

\item Supongamos ahora que $A$ y $B$ son \emph{normales} y conmutan. Usando (a) y el
Teorema~\ref{teo:normal-iff-ortonormal} aplicado a la restricción de $B$ a cada
autoespacio de $A$, construir una base ortonormal de $\mathbb{C}^{n}$ formada por
autovectores \emph{de las dos matrices a la vez}. Concluir que existe una única unitaria
$U$ con
\[
U^{*}AU=D_1 \quad\text{y}\quad U^{*}BU=D_2
\]
ambas diagonales. \emph{Sugerencia:} hay que verificar que la restricción de $B$ a un
autoespacio de $A$ sigue siendo normal.

\item Probar el recíproco: si existe una unitaria que diagonaliza a las dos, entonces
conmutan. \emph{(Es la parte fácil: las diagonales conmutan.)}

\item Verificar que dos matrices circulantes cualesquiera conmutan (Ejercicio~5 de este
capítulo) y explicar, a la luz de (b) y (c), por qué \emph{todas} las circulantes de
tamaño $n$ se diagonalizan con la misma matriz de Fourier $F$.

\item Dar dos matrices normales que no conmuten y comprobar que no tienen ninguna base
ortonormal de autovectores en común.

\item \emph{(Una consecuencia.)} Probar que si $A$ es normal y $p$ es un polinomio,
entonces $A$ y $p(A)$ conmutan, y deducir que se diagonalizan simultáneamente con los
mismos autovectores y autovalores $\lambda_k$ y $p(\lambda_k)$.
\end{enumerate}

\medskip
\noindent{\footnotesize\emph{Se entrega por escrito. El ítem (b) es la razón profunda por
la que en procesamiento de señales una sola transformada, la de Fourier, sirve para
todos los filtros a la vez.}}
\end{entrega}

\section{Los casos particulares importantes}\label{sec:casos-particulares}

\begin{corollary}[Caso hermítico]\label{cor:hermitico}
Si \(A\) es hermítica (\(A^{*}=A\)), entonces \(A\) es normal y, por el
Teorema~\ref{teo:normal-iff-ortonormal}, existen $U$ unitaria y $\Lambda$ diagonal
\emph{real} tales que \(A=U\Lambda U^{*}\).
En particular, los autovalores de una matriz hermítica son \emph{reales}.
\end{corollary}

\begin{proof}
Si \(A^{*}=A\), entonces \(AA^{*}=A^{2}=A^{*}A\): \(A\) es normal.
La diagonalización unitaria sigue del Teorema~\ref{teo:normal-iff-ortonormal}.
Que $\Lambda$ sea real se deduce de $Av=\lambda v$ con $\|v\|=1$:
\[
\lambda= \lambda\|v\|^{2}=v^{*}Av=v^{*}A^{*}v=\overline{v^{*}Av}=\overline{\lambda}
\;\Longrightarrow\;\lambda\in\mathbb{R}.
\]
\end{proof}

\begin{corollary}[Antihermíticas y unitarias]   % <-- NUEVO
\begin{enumerate}[label=(\roman*)]
\item Si $A^{*}=-A$ (antihermítica), sus autovalores son \emph{imaginarios puros}.
\item Si $A^{*}A=I$ (unitaria), sus autovalores tienen \emph{módulo $1$}.
\end{enumerate}
En ambos casos $A$ es normal, luego unitariamente diagonalizable.
\end{corollary}

\begin{proof}
(i) $A$ es normal ($AA^{*}=-A^{2}=A^{*}A$) y, si $Av=\lambda v$ con $\|v\|=1$,
$\lambda=v^{*}Av=-v^{*}A^{*}v=-\overline{\lambda}$, o sea $\operatorname{Re}\lambda=0$.

\smallskip
(ii) $A$ es normal pues $A^{*}=A^{-1}$ conmuta con $A$. Si $Av=\lambda v$ con $\|v\|=1$,
entonces $1=\|v\|^2=\|Av\|^{2}=|\lambda|^{2}$.
\end{proof}

\paragraph{Particularización al caso real.}   % <-- NUEVO
Si $A\in\mathbb{R}^{n\times n}$, la adjunta es la traspuesta y ``normal'' significa
$AA^{T}=A^{T}A$. Hay que tener cuidado: \emph{una matriz real normal puede no tener
autovalores reales}, y en ese caso no hay diagonalización real (ver
Ejemplo~\ref{ej:rotacion}). El caso perfecto es el simétrico:

\begin{theorem}[Teorema espectral real]\label{teo:espectral-real}   % <-- NUEVO
Si $A\in\mathbb{R}^{n\times n}$ es simétrica ($A^{T}=A$), existen $Q$ \emph{ortogonal}
real ($Q^{T}Q=I$) y $\Lambda$ diagonal real tales que
\[
A=Q\Lambda Q^{T}=\sum_{k=1}^{n}\lambda_k\,q_kq_k^{T}.
\]
\end{theorem}

\begin{proof}
Vista como matriz compleja, $A$ es hermítica, así que por el
Corolario~\ref{cor:hermitico} todos sus autovalores $\lambda$ son reales.
Para cada uno, el sistema $(A-\lambda I)x=0$ tiene coeficientes reales, luego admite
una base de soluciones \emph{reales}; aplicando Gram--Schmidt dentro de cada autoespacio
obtenemos bases ortonormales reales. Por el Lema~\ref{lema:normal-basico}(iii),
autoespacios de autovalores distintos son ortogonales entre sí, y por el
Teorema~\ref{teo:normal-iff-ortonormal} la suma de sus dimensiones es $n$.
Juntando todas esas bases se obtiene una base ortonormal real de autovectores,
que forma las columnas de $Q$.
\end{proof}

\begin{example}[Una rotación: normal real sin autovectores reales]\label{ej:rotacion}  % <-- NUEVO
Sea
\[
R_\theta=\begin{bmatrix}\cos\theta & -\sen\theta\\ \sen\theta & \cos\theta\end{bmatrix},
\qquad 0<\theta<\pi .
\]
Es ortogonal, luego normal. Su polinomio característico es $z^{2}-2\cos\theta\,z+1$,
con raíces $e^{\pm i\theta}$: son de módulo $1$ (como debe ser) pero \emph{no reales}.
Geométricamente es evidente: una rotación no deja ninguna recta de $\mathbb{R}^{2}$ fija.
Sobre $\mathbb{C}^{2}$, en cambio, $R_\theta$ es unitariamente diagonalizable con
autovectores $\tfrac1{\sqrt2}(1,\mp i)^{T}$.
\end{example}

\begin{theorem}[Forma de Schur real]\label{teo:schur-real}
Si $A\in\mathbb{R}^{n\times n}$, existe $Q$ ortogonal \emph{real} tal que $Q^{T}AQ$ es
triangular superior \emph{por bloques}, con bloques diagonales de tamaño $1\times1$ (los
autovalores reales) y $2\times2$ (uno por cada par de autovalores complejos conjugados,
como en el Ejemplo~\ref{ej:rotacion}). Es la forma que efectivamente calculan los
algoritmos numéricos, porque se obtiene sin salir de la aritmética real.
\end{theorem}

La demostración es el ejercicio de desarrollo que sigue: es la misma inducción del
Teorema~\ref{teo:schur}, con el cuidado de no abandonar $\mathbb{R}$.

\begin{entrega}{La forma de Schur real}\label{ent:schur-real}
Sea $A\in\mathbb{R}^{n\times n}$.

\begin{enumerate}[label=\textbf{(\alph*)},itemsep=3pt,topsep=4pt]
\item Supongamos que $A$ tiene un autovalor \emph{real} $\lambda$. Probar que existe un
autovector real asociado. \emph{Sugerencia:} el sistema $(A-\lambda I)x=0$ tiene
coeficientes reales.

\item Supongamos ahora $Av=\lambda v$ con $\lambda=\alpha+i\beta$, $\beta\neq0$, y
escribamos $v=p+iq$ con $p,q\in\mathbb{R}^{n}$. Separando parte real e imaginaria,
probar que
\[
Ap=\alpha p-\beta q,
\qquad
Aq=\beta p+\alpha q .
\]
Deducir que $W=\gen\{p,q\}$ es un subespacio \emph{real} invariante por $A$.

\item Probar que $p$ y $q$ son linealmente independientes, de modo que $\dim W=2$.
\emph{Sugerencia:} si $q=cp$ entonces $v=(1+ic)p$ sería un múltiplo complejo de un vector
real, y entonces $Av=\lambda v$ obligaría a $\lambda$ a ser real.

\item Ortonormalizar $\{p,q\}$ con Gram--Schmidt para obtener $q_1,q_2$, verificar que la
matriz de $A$ restringida a $W$ en esa base es de $2\times2$ y comprobar que sus
autovalores son $\lambda$ y $\overline{\lambda}$.

\item Armar la inducción: completar $q_1,q_2$ a una base ortonormal de $\mathbb{R}^{n}$,
escribir $Q_1^{T}AQ_1$ por bloques y aplicar la hipótesis inductiva al bloque de abajo a
la derecha. Tratar por separado el caso de autovalor real, donde el bloque es $1\times1$.

\item Hacer las cuentas explícitamente con
\[
A=\begin{bmatrix} 1&0&0\\ 0&0&-1\\ 0&1&0\end{bmatrix},
\]
identificando el bloque $1\times1$ y el bloque $2\times2$.

\item Explicar por qué un algoritmo que trabaja en aritmética real no puede producir la
forma triangular compleja del Teorema~\ref{teo:schur}, y qué se gana a cambio con esta
versión.
\end{enumerate}

\medskip
\noindent{\footnotesize\emph{Se entrega por escrito. El ítem (b) es el mecanismo por el
cual un par de autovalores complejos conjugados de una matriz real se manifiesta como una
\emph{rotación} en un plano, que es lo que vimos en el Capítulo~\ref{cap:autovalores} al
resolver $\n{u}'=A\n{u}$.}}
\end{entrega}

\begin{entrega}{Caracterizaciones de las matrices normales}\label{ent:normales}
Sea $A\in\mathbb{C}^{n\times n}$.

\begin{enumerate}[label=\textbf{(\alph*)},itemsep=3pt,topsep=4pt]
\item Probar que $A$ es normal si y sólo si $\|A\n{x}\|=\|A^{*}\n{x}\|$ para todo
$\n{x}$. \emph{Sugerencia:} una implicación es el Lema~\ref{lema:normal-basico}(i);
para la otra, desarrollar $\n{x}^{*}(A^{*}A-AA^{*})\n{x}=0$ y usar que una matriz
hermítica con esa propiedad es nula (polarización).

\item Deducir que si $A$ es normal entonces $N(A)=N(A^{*})$ y, tomando complementos
ortogonales, $C(A)=C(A^{*})$. Concluir que
\[
\mathbb{C}^{n}=C(A)\oplus N(A)
\]
es una descomposición \emph{ortogonal}: para una matriz normal, los cuatro subespacios
fundamentales del Capítulo~\ref{cap:ortogonalidad} colapsan a dos.

\item Dar una matriz con $N(A)=N(A^{*})$ que no sea normal, para ver que (b) no se
invierte.

\item Probar que $A$ es normal si y sólo si los autovalores de $A^{*}A$ son exactamente
$|\lambda_1|^{2},\dots,|\lambda_n|^{2}$, donde los $\lambda_i$ son los autovalores de
$A$. \emph{Sugerencia:} comparar trazas usando la desigualdad de
Schur~\ref{teo:desig-schur}.

\item \emph{(Una caracterización sorprendente.)} Probar que si $A$ es normal entonces
existe un polinomio $p$ con $A^{*}=p(A)$. \emph{Sugerencia:} diagonalizar y buscar un
polinomio que interpole $p(\lambda_k)=\overline{\lambda_k}$ en los autovalores distintos.
Probar también el recíproco.

\item Verificar (e) explícitamente para una matriz de permutación de $3\times3$ y para
una simétrica real (¿cuál es $p$ en ese caso?).
\end{enumerate}

\medskip
\noindent{\footnotesize\emph{Se entrega por escrito. El ítem (b) es la razón por la que
para matrices normales --- y en particular simétricas --- no hace falta distinguir entre
el espacio fila y el espacio columna.}}
\end{entrega}

\section{Una consecuencia cuantitativa: cuánto se ``ve'' en los autovalores}\label{sec:desig-schur}

\begin{theorem}[Desigualdad de Schur]\label{teo:desig-schur}
Sea $A\in\mathbb{C}^{n\times n}$ con autovalores $\lambda_1,\dots,\lambda_n$
(con multiplicidad). Entonces
\[
\sum_{k=1}^{n}|\lambda_k|^{2}\;\le\;\|A\|_F^{2}:=\sum_{i,j}|a_{ij}|^{2},
\]
con igualdad \emph{si y solo si} $A$ es normal.
\end{theorem}

\begin{proof}
La norma de Frobenius, definida en la Sección~\ref{sec:frobenius}, satisface
$\|A\|_F^{2}=\tr(A^{*}A)$ y es invariante por multiplicación por unitarias: si
$U^{*}U=I$,
\[
\|UA\|_F^{2}=\tr\big((UA)^{*}UA\big)=\tr(A^{*}U^{*}UA)=\tr(A^{*}A)=\|A\|_F^{2},
\]
y análogamente por la derecha usando $\tr(XY)=\tr(YX)$; en particular
$\|UAV\|_F=\|A\|_F$. Si $A=UTU^{*}$ es la
descomposición de Schur, entonces
\[
\|A\|_F^{2}=\|T\|_F^{2}=\sum_{k}|\lambda_k|^{2}+\sum_{i<j}|t_{ij}|^{2}
\;\ge\;\sum_{k}|\lambda_k|^{2},
\]
con igualdad si y solo si todos los $t_{ij}$ con $i<j$ son nulos, es decir,
si y solo si $T$ es diagonal, lo que por el Teorema~\ref{teo:normal-iff-ortonormal}
equivale a que $A$ sea normal.
\end{proof}

\begin{remark}
Para una matriz normal, entonces, ``toda la energía'' de $A$ está en sus autovalores:
$\|A\|_F^{2}=\sum_k|\lambda_k|^{2}$. Esto es lo que permite medir el error de
aproximar $A$ quedándose solo con los autovalores más grandes, que es el contenido
del Ejercicio de desarrollo~\ref{ent:rango-bajo}.
La diferencia $\|A\|_F^{2}-\sum_k|\lambda_k|^{2}$ mide, en cambio, cuán lejos de
normal está $A$.
\end{remark}

\paragraph{Más ejemplos de matrices normales.}
\begin{itemize}
\item \textbf{Unitarias:} \(U^{*}U=I\Rightarrow UU^{*}=I\), luego \(UU^{*}=U^{*}U\).
\item \textbf{Ortogonales} (caso real de las unitarias); en particular, las
      \textbf{matrices de permutación}.   % <-- NUEVO (se usa en el Ejercicio 5)
\item \textbf{Hermíticas y simétricas reales.}
\item \textbf{Skew-hermíticas (antihermíticas):} \(A^{*}=-A\Rightarrow AA^{*}=-A^{2}=A^{*}A\).
\item \textbf{Diagonales} (y, por semejanza unitaria, todas las unitariamente
      diagonalizables): conmutan con su adjunta.
\item \textbf{Bloque-diagonales:} \(\operatorname{diag}(N_{1},\dots,N_{k})\) es normal
      si y solo si cada \(N_{i}\) es normal.
\item \textbf{Polinomios de una normal:} si \(A\) es normal y \(p\) es un polinomio,
      entonces \(p(A)\) es normal (de \(A=UDU^{*}\) sale \(p(A)=Up(D)U^{*}\)).
\end{itemize}

\begin{entrega}{Aproximación de rango bajo de una matriz simétrica}
\label{ent:rango-bajo}
Sea $S\in\mathbb{R}^{n\times n}$ simétrica, con autovalores ordenados por módulo
$|\lambda_1|\ge|\lambda_2|\ge\cdots\ge|\lambda_n|$ y base ortonormal de autovectores
$q_1,\dots,q_n$, de modo que $S=\sum_{k}\lambda_kq_kq_k^{T}$
(Teorema~\ref{teo:espectral-real}). Para $1\le r<n$ definimos el \emph{truncamiento}
\[
S_r=\sum_{k=1}^{r}\lambda_k\,q_kq_k^{T}.
\]

\begin{enumerate}[label=\textbf{(\alph*)},itemsep=3pt,topsep=4pt]
\item Probar que $S_r$ es simétrica, que $\rg(S_r)\le r$, y que
\[
\|S-S_r\|_F^{2}=\lambda_{r+1}^{2}+\cdots+\lambda_n^{2}.
\]

\item Probar que $\|S\|_F^{2}=\sum_k\lambda_k^{2}$ (es la
igualdad del Teorema~\ref{teo:desig-schur} para matrices normales) y deducir que el
error relativo de la aproximación es
$\|S-S_r\|_F^{2}/\|S\|_F^{2}=\sum_{k>r}\lambda_k^{2}\big/\sum_k\lambda_k^{2}$.

\item \emph{(Optimalidad.)} Probar que $S_r$ es la \emph{mejor} aproximación: para toda
matriz simétrica $B$ con $\rg(B)\le r$,
\[
\|S-B\|_F\ \ge\ \|S-S_r\|_F .
\]
\emph{Sugerencia:} escribir todo en la base $q_1,\dots,q_n$, es decir, reemplazar $S$ y
$B$ por $Q^{T}SQ=\Lambda$ y $Q^{T}BQ$, lo cual no cambia ninguna de las dos normas de
Frobenius ni el rango; después mirar la diagonal.

\item Aplicarlo a la matriz de similitudes del Ejercicio~4 de este capítulo y comparar
con el cálculo directo que se pide allí.

\item ¿Qué pasa si se ordenan los autovalores por valor y no por módulo? Dar un ejemplo
de $2\times2$ donde eso produzca una aproximación peor.

\end{enumerate}

\medskip
\noindent{\footnotesize\emph{Se entrega por escrito. El ítem (c) es la primera versión
que vemos de un teorema de optimalidad: truncar no sólo da \emph{una} aproximación de
rango bajo, da la mejor que existe.}}
\end{entrega}

\begin{observacion}[Hacia dónde va esto]\label{obs:hacia-rango-bajo}
El ejercicio anterior es un anticipo de dos resultados que se tratan más adelante y que
conviene no mezclar con lo de acá.

El primero es la \textbf{versión general}. Lo que probamos vale para matrices
simétricas, que son las que tienen base ortonormal de autovectores. Para una matriz
cualquiera --- incluso rectangular --- el enunciado análogo es el \emph{teorema de
Eckart-Young}, donde los autovalores se reemplazan por los valores singulares y el
truncamiento $S_r$ por el de la descomposición en valores singulares. La demostración
sigue las mismas líneas. Los valores singulares aparecen en el
Capítulo~\ref{cap:SVD}.

El segundo es la \textbf{lectura estadística}. Si las filas de una matriz $X$ son
observaciones de varias variables y uno le resta a cada columna su media, la matriz
$X^{T}X$ es simétrica y semidefinida positiva, y sus autovectores resultan ser las
direcciones en las que los datos más se dispersan. Truncar como en este ejercicio es
entonces quedarse con esas pocas direcciones y describir la nube de puntos con menos
coordenadas de las que tenía. Esa técnica se llama \emph{análisis de componentes
principales}, y tiene capítulo propio después de la SVD; recién ahí estaremos en
condiciones de decir en qué sentido preciso esas direcciones son las mejores.
\end{observacion}

\section{Normal, diagonalizable, defectiva}\label{sec:panorama}

Conviene ordenar las tres nociones que estuvimos manejando, porque es fácil confundirlas.
Para $A\in\mathbb{C}^{n\times n}$:

\begin{center}
\renewcommand{\arraystretch}{1.35}
\begin{tabular}{@{}lll@{}}
\hline
 & \textbf{qué existe} & \textbf{caracterización}\\
\hline
siempre & base ortonormal con $A$ \emph{triangular} & Teorema~\ref{teo:schur}\\
diagonalizable & base de autovectores \emph{independientes}
   & mult.\ geom.\ $=$ mult.\ alg.\\
normal & base \emph{ortonormal} de autovectores
   & $AA^{*}=A^{*}A$ (Teorema~\ref{teo:normal-iff-ortonormal})\\
\hline
\end{tabular}
\end{center}

Las tres filas son cada vez más exigentes, y ninguna de las dos implicaciones se
invierte. Que ser diagonalizable no alcance para ser normal se ve en el
Ejemplo~\ref{ej:schur-numerico}: la matriz
$\begin{bsmallmatrix}1&1\\ 0&2\end{bsmallmatrix}$ tiene autovalores distintos, luego es
diagonalizable por el criterio del Capítulo~\ref{cap:autovalores}, pero no es normal y
sus autovectores $(1,0)^{T}$ y $(1,1)^{T}$ son independientes pero no ortogonales. Y que
Schur no alcance para ser diagonalizable se ve en el bloque de Jordan
$\begin{bsmallmatrix}\lambda&1\\ 0&\lambda\end{bsmallmatrix}$, que es \emph{defectiva}:
tiene un solo autovector independiente. Schur, en cambio, no falla nunca; en ese caso lo
mejor que se puede conseguir es, justamente, triangular.

\begin{remark}[Lo que viene]
El Corolario~\ref{cor:espectral} escribe una matriz normal como
$A=\sum_k\lambda_k\,q_kq_k^{*}$: una suma de proyecciones ortogonales pesadas por los
autovalores, con los $q_k$ ortonormales. En el Capítulo~\ref{cap:SVD} veremos que
\emph{toda} matriz, incluso rectangular, admite una descomposición del mismo aspecto,
\[
A=\sum_{k=1}^{r}\sigma_k\,u_k v_k^{T},
\]
con dos familias ortonormales en lugar de una y con los $\sigma_k>0$. El precio de
renunciar a la normalidad --- y hasta a que la matriz sea cuadrada --- es tener que usar
dos bases distintas en lugar de una. La herramienta para construirla es el
Teorema~\ref{teo:espectral-real} aplicado a la matriz simétrica $A^{T}A$.
\end{remark}

% ==========================================================
\section{Ejercicios}
\begin{enumerate}

\item {\bf Matrices de covarianza, traza y ortogonalidad}\\
En Ciencia de Datos es común trabajar con la matriz de covarianza (sin normalizar)
$S = X^{T}X$, donde $X \in \mathbb{R}^{n \times p}$ es una matriz de datos.
\begin{enumerate}[label=\alph*)]
    \item Demostrá que $S$ es simétrica y, en consecuencia, normal.
    \item Usando el Teorema~\ref{teo:espectral-real}, justificá por qué $S$ admite una
          matriz de autovectores $Q$ con $Q^{-1}=Q^{T}$ (matriz ortogonal),
          y por qué sus autovalores son reales.
    \item Sabiendo que la traza de una matriz es igual a la suma de sus autovalores,
          demostrá que la suma de los autovalores de $S$ es igual a la suma de las
          normas al cuadrado de las columnas de $X$.
    \item Probá que $\lambda_k=\|Xq_k\|^{2}\ \ge 0$ para todo $k$, donde $q_k$ es el
          $k$-ésimo autovector unitario de $S$. Interpretá: ¿qué mide $\lambda_k$ sobre
          los datos proyectados en la dirección $q_k$?   % <-- NUEVO
\end{enumerate}

\item {\bf Proyecciones ortogonales y normalidad}\\
Sea $Q \in \mathbb{R}^{n \times k}$ una matriz alta y delgada con columnas ortonormales,
es decir, $Q^{T}Q = I_k$. Definimos $P = QQ^{T}$.
\begin{enumerate}[label=\alph*)]
    \item Demostrá que $P$ es normal.
    \item Verificá analíticamente que $P^{2}=P$.
    \item Usando la definición de autovalor ($Pv=\lambda v$, $v\neq0$) y la propiedad
          anterior, demostrá que los únicos autovalores posibles de $P$ son
          $\lambda=0$ y $\lambda=1$.
    \item ¿Es $P$ diagonalizable ortogonalmente? Justificá con el
          Teorema~\ref{teo:normal-iff-ortonormal}.
    \item Deducí, usando la traza, que la multiplicidad del autovalor $1$ es exactamente
          $k$. ¿Cuál es el autoespacio asociado a $\lambda=1$? ¿Y a $\lambda=0$?  % <-- NUEVO
\end{enumerate}

\item {\bf Matrices de adyacencia y análisis de redes}\\
Un grafo no dirigido se representa por su matriz de adyacencia
$W \in \mathbb{R}^{n \times n}$, con $W_{ij}=W_{ji}\ge 0$ el peso de la conexión
entre los nodos $i$ y $j$.
\begin{enumerate}[label=\alph*)]
    \item Justificá que $W$ es simétrica real, luego normal, y que por lo tanto
          sus autovalores son reales.   % <-- CORREGIDO: antes decía solo "por qué es normal"
    \item Sea $D$ diagonal con $D_{ii}=\sum_{j=1}^{n}W_{ij}$ (el grado del nodo $i$).
          Definimos la \emph{matriz Laplaciana} $L=D-W$. Demostrá que $L$ es normal.
    \item Demostrá que el vector de unos $\mathbf{1}=(1,\dots,1)^{T}$ es siempre
          autovector de $L$. ¿Cuál es el autovalor asociado?
    \item Probá la \emph{identidad de Dirichlet}: para todo $x\in\mathbb{R}^{n}$,   % <-- NUEVO
          \[
            x^{T}Lx=\tfrac12\sum_{i,j=1}^{n}W_{ij}\,(x_i-x_j)^{2}\;\ge\;0 .
          \]
          Deducí que todos los autovalores de $L$ son $\ge 0$.
    \item Supongamos que el grafo tiene dos componentes conexas $C_1$ y $C_2$.
          Exhibí dos autovectores linealmente independientes de $L$ con autovalor $0$
          (sugerencia: vectores indicadores de $C_1$ y $C_2$, y usá (d) para ver que
          $Lx=0 \iff x$ es constante en cada componente). ¿Qué sugiere esto como
          algoritmo para \emph{detectar} comunidades en una red?   % <-- NUEVO
\end{enumerate}

\item {\bf Descomposición espectral y aproximación de rango bajo}   % <-- NUEVO
\\
Sea
\[
A=\begin{pmatrix}2&1&1\\1&2&1\\1&1&2\end{pmatrix},
\]
la matriz de similitudes entre tres documentos.
\begin{enumerate}[label=\alph*)]
    \item Escribí $A=I+\mathbf{1}\mathbf{1}^{T}$ y hallá, sin calcular el polinomio
          característico, todos los autovalores de $A$ y una base ortonormal de autovectores.
          (Sugerencia: ¿qué le hace $\mathbf{1}\mathbf{1}^{T}$ a $\mathbf 1$? ¿y a un
          vector ortogonal a $\mathbf 1$?)
    \item Escribí explícitamente la descomposición espectral
          $A=\sum_{k=1}^{3}\lambda_kq_kq_k^{T}$ del Corolario~\ref{cor:espectral}.
    \item Sea $A_1=\lambda_1q_1q_1^{T}$ con $\lambda_1$ el autovalor de mayor módulo.
          Calculá la matriz $A-A_1$ y verificá que
          $\|A-A_1\|_F=\sqrt{\lambda_2^{2}+\lambda_3^{2}}$, como predice el
          Teorema~\ref{teo:desig-schur}.
    \item Interpretá $q_1$: ¿qué característica de los tres documentos captura esa
          dirección? ¿Y qué información se pierde al quedarse solo con $A_1$?
    \item Calculá $A^{m}$ para todo $m$ usando el Corolario~\ref{cor:espectral}, y
          mostrá que $\dfrac{A^{m}}{\lambda_1^{m}}\longrightarrow q_1q_1^{T}$
          cuando $m\to\infty$. (Esta es la idea del \emph{método de la potencia}.)
\end{enumerate}

\item {\bf Filtros, matrices circulantes y Fourier}   % <-- NUEVO
\\
Sea $P\in\mathbb{R}^{n\times n}$ la permutación cíclica definida por
$Pe_j=e_{j+1}$ para $j=1,\dots,n$ (con $e_{n+1}:=e_1$); en coordenadas,
$(Px)_i=x_{i-1}$ (con $x_0:=x_n$). Una \emph{matriz circulante} es
\[
C=c_0I+c_1P+c_2P^{2}+\dots+c_{n-1}P^{\,n-1},\qquad c_j\in\mathbb{R}.
\]
\begin{enumerate}[label=\alph*)]
    \item Probá que $P$ es ortogonal ($P^{T}P=I$) y deducí que es normal.
          ¿Es simétrica? ¿Cuánto vale $P^{n}$?
    \item Deducí que $C$ es normal para cualquier elección de los $c_j$.
    \item Sea $\omega=e^{2\pi i/n}$ y
          $f_k=\tfrac{1}{\sqrt n}\bigl(1,\omega^{k},\omega^{2k},\dots,\omega^{(n-1)k}\bigr)^{T}$,
          $k=0,\dots,n-1$. Verificá que $Pf_k=\omega^{-k}f_k$ y que
          $\{f_0,\dots,f_{n-1}\}$ es una base ortonormal de $\mathbb{C}^{n}$
          (la matriz $F=[f_0\;\cdots\;f_{n-1}]$ es la \emph{matriz de Fourier}, unitaria).
    \item Concluí que \emph{todas} las circulantes de tamaño $n$ se diagonalizan con la
          \emph{misma} unitaria $F$, con autovalores
          \[
            \lambda_k=\sum_{j=0}^{n-1}c_j\,\omega^{-jk},
          \]
          es decir, la transformada discreta de Fourier del vector de coeficientes $c$.
    \item Deducí que dos matrices circulantes cualesquiera \emph{conmutan}, y que el
          producto de circulantes es circulante. (En el lenguaje de señales:
          convolucionar es multiplicar coordenada a coordenada en el dominio de Fourier.)
    \item {\bf Suavizado.} Considerá el filtro de \emph{media móvil} de tres términos:
          $c_0=c_1=c_{n-1}=\tfrac13$ y el resto nulo, de modo que
          $(Cx)_i=\tfrac13(x_{i-1}+x_i+x_{i+1})$. Mostrá que
          \[
            \lambda_k=\tfrac13\Bigl(1+2\cos\tfrac{2\pi k}{n}\Bigr).
          \]
          Interpretá: ¿qué le pasa a la componente de frecuencia baja ($k=0$)?
          ¿Y a la de frecuencia alta ($k\approx n/2$)?
    \item Usando el Corolario~\ref{cor:espectral}, describí qué le pasa a una señal
          $x$ al aplicarle el filtro $m$ veces, es decir, calculá $C^{m}x$ en términos
          de los coeficientes de Fourier de $x$. ¿Hacia qué señal converge $C^{m}x$
          cuando $m\to\infty$?
\end{enumerate}

\item {\bf Falso o verdadero} (justificar o dar contraejemplo)   % <-- NUEVO
\begin{enumerate}[label=\alph*)]
    \item Si $A$ es normal y $\lambda$ es autovalor de $A$, entonces $\overline{\lambda}$
          es autovalor de $A$.
    \item Si $A$ y $B$ son normales, $A+B$ es normal.
    \item Si $A$ es normal e invertible, $A^{-1}$ es normal.
    \item Si $A$ es real y normal, entonces $A$ es simétrica.
    \item Si $A$ es triangular superior y todos sus autovalores son reales, entonces
          $A$ es diagonalizable.
    \item Si $\|A\|_F^{2}=\sum_k|\lambda_k|^{2}$, entonces $A$ admite una base
          ortonormal de autovectores.
\end{enumerate}

\end{enumerate}